% ============================================================================
% Fichier  : partie4/P04_C17_identites_acces.tex
% Titre    : Identités et Accès
% Auteur   : MINKA MI NGUIDJOÏ Thierry Emmanuel
% Version  : 1.0 -- novembre 2026
% ============================================================================
% Base : brouillon « Livre VI » (IAAA, authentification, contrôle d'accès).
% Sources : NIST SP 800-63B (2017) ; W3C Web Authentication (2019) ;
% OASIS SAML 2.0 (2005) ; OpenID Connect Core 1.0 (2014) ; CIS Controls 5
% et 6. Chiffres de la revue des comptes : cas fictif, cohérents entre eux.
% ============================================================================

\chapter{Identités et Accès}
\label{chap:iam}
\chaptermark{Identités et accès}

\epigraph{%
    « Un mot de passe prouve que l'on connaît un secret,\\
    pas que l'on est la bonne personne. »%
}{Principe du cours}

\niveauChapitre{Intermédiaire}{Chapitres~\ref{chap:crypto-sym},
\ref{chap:crypto-asym} et~\ref{chap:modeles}.}

\begin{boxobjectifs}
\begin{itemize}
    \item Distinguer identification, authentification, autorisation et
          traçabilité.
    \item Choisir des facteurs d'authentification et une politique de mots
          de passe fondée sur le NIST.
    \item Expliquer l'authentification résistante à l'hameçonnage.
    \item Décrire l'authentification unique (SAML, OpenID Connect).
    \item Gérer le cycle de vie des comptes et les comptes à privilèges.
    \item Conduire une revue des droits.
\end{itemize}
\end{boxobjectifs}

\bigskip

% ============================================================================
\section*{Question directrice : qui pouvait se connecter au système de LiZbiZ~?}
% ============================================================================

Après l'attaque, LiZbiZ extrait la liste des comptes de son annuaire~: 212
comptes actifs pour 150~salariés. Parmi eux, des comptes de salariés
partis, des comptes partagés, des comptes de service jamais revus, et le
compte du prestataire, utilisé par plusieurs techniciens. Personne ne
pouvait dire qui avait réellement accès à quoi. C'est le sujet de ce
chapitre.

\medskip
\textbf{Plan du chapitre.}
\begin{enumerate}
    \item Identifier, authentifier, autoriser, tracer (section~\ref{sec:c17-iaaa}).
    \item Les facteurs et les mots de passe (section~\ref{sec:c17-facteurs}).
    \item L'authentification forte (section~\ref{sec:c17-mfa}).
    \item L'authentification unique et les annuaires (section~\ref{sec:c17-sso}).
    \item Le cycle de vie et les privilèges (section~\ref{sec:c17-cycle}).
    \item La revue des droits (section~\ref{sec:c17-revue}).
\end{enumerate}

\separateur

% ============================================================================
\section{Identifier, Authentifier, Autoriser, Tracer}
\label{sec:c17-iaaa}
% ============================================================================

\begin{tableaucours}{L{3.2cm}L{4.6cm}L{4.4cm}}{%
    \tableauHeader{Étape} &
    \tableauHeader{Question} &
    \tableauHeader{Exemple}
}
Identification & Qui prétendez-vous être~? & Saisir son identifiant \\
Authentification & Le prouvez-vous~? & Mot de passe et second facteur \\
Autorisation & Qu'avez-vous le droit de faire~? & Le rôle « comptable »
donne accès aux paiements \\
Traçabilité & Qu'avez-vous fait~? & Journal des actions, compte
nominatif \\
\end{tableaucours}
\vspace{-0.8em}
\begin{center}{\small\itshape Tableau~17.1. Le cadre IAAA}\end{center}

\begin{boxjuris}
\textbf{\loicam{n°~2010/012}, article~4, point~10} (texte officiel)~:
\textit{« Authentification~: critère de sécurité défini par un processus
mis en oeuvre notamment pour vérifier l'identité d'une personne physique
ou morale et s'assurer que l'identité fournie correspond à l'identité de
cette personne préalablement enregistrée. »}
\end{boxjuris}

Un compte partagé ruine la traçabilité~: si cinq techniciens utilisent le
même compte, aucun journal ne dira lequel a agi.

% ============================================================================
\section{Les Facteurs et les Mots de Passe}
\label{sec:c17-facteurs}
% ============================================================================

On distingue trois catégories de facteurs~: ce que l'on \emph{sait}
(mot de passe, code), ce que l'on \emph{possède} (téléphone, clé
physique, carte) et ce que l'on \emph{est} (empreinte, visage).
L'authentification \emph{multifacteur} combine des facteurs de catégories
différentes~: deux mots de passe ne font pas deux facteurs.

\begin{tableaucours}{L{5.0cm}L{7.2cm}}{%
    \tableauHeader{Recommandation du NIST SP~800-63B} &
    \tableauHeader{Raison}
}
Longueur minimale, et accepter les longues phrases & La longueur fait la
résistance~\cite{nist80063b} \\
Pas de règles de composition imposées & Elles produisent des mots de passe
prévisibles (« Lizbiz2026! ») \\
Pas de changement périodique imposé & On change seulement en cas de
compromission~; sinon les utilisateurs incrémentent \\
Comparer aux listes de mots de passe divulgués & Bloquer les mots de passe
déjà connus des attaquants \\
Limiter les tentatives & Contre la force brute en ligne \\
\end{tableaucours}
\vspace{-0.8em}
\begin{center}{\small\itshape Tableau~17.2. Politique de mots de passe
fondée sur le NIST}\end{center}

% ============================================================================
\section{L'Authentification Forte}
\label{sec:c17-mfa}
% ============================================================================

Tous les seconds facteurs ne se valent pas. Un code reçu par SMS peut être
détourné (échange de carte SIM, Chapitre~\ref{chap:menaces})~; un code
saisi sur une fausse page peut être relayé en temps réel par l'attaquant.
Les authentificateurs \emph{résistants à l'hameçonnage} (clés FIDO2,
WebAuthn) signent un défi lié au domaine du site~: sur une fausse page, ils
ne répondent pas~\cite{w3c2019webauthn}.

\begin{tableaucours}{L{3.6cm}L{4.4cm}L{4.2cm}}{%
    \tableauHeader{Second facteur} &
    \tableauHeader{Forces} &
    \tableauHeader{Faiblesses}
}
Code par SMS & Simple, sans équipement & Échange de SIM, interception,
relais \\
Application à code (TOTP) & Hors du réseau téléphonique & Relayable sur
une fausse page \\
Notification à valider & Simple & Lassitude~: l'utilisateur valide une
demande qu'il n'a pas faite \\
Clé FIDO2, WebAuthn & Résiste à l'hameçonnage & Coût, gestion des pertes \\
\end{tableaucours}
\vspace{-0.8em}
\begin{center}{\small\itshape Tableau~17.3. Les seconds facteurs}\end{center}

% ============================================================================
\section{L'Authentification Unique et les Annuaires}
\label{sec:c17-sso}
% ============================================================================

Un \termecle{annuaire} (Active Directory, LDAP) centralise les comptes et
les groupes. L'\termecle{authentification unique} (SSO) permet de
s'authentifier une fois auprès d'un fournisseur d'identité, qui atteste
ensuite l'identité auprès des applications~: par des assertions SAML~2.0
\cite{oasis2005saml} ou des jetons OpenID Connect~\cite{oidc2014}. Le SSO
réduit le nombre de mots de passe et centralise la double
authentification~; il concentre aussi le risque~: le fournisseur
d'identité devient une cible majeure, à protéger en priorité.

% ============================================================================
\section{Le Cycle de Vie et les Privilèges}
\label{sec:c17-cycle}
% ============================================================================

\begin{boxprocedure}
\textbf{Le cycle de vie d'un compte}
\begin{enumerate}
    \item \textbf{Arrivée}~: création sur demande des ressources humaines,
          avec le rôle validé par le responsable.
    \item \textbf{Mobilité}~: retrait des droits de l'ancien poste avant
          l'ajout de ceux du nouveau.
    \item \textbf{Absence longue}~: suspension.
    \item \textbf{Départ}~: désactivation le jour même, suppression après
          le délai fixé.
    \item \textbf{Revue périodique}~: rapprochement des comptes avec le
          personnel et des droits avec les rôles.
\end{enumerate}
\end{boxprocedure}

\begin{boxdef}
\termecle{Compte à privilèges}~: compte qui peut modifier la configuration,
les droits ou la sécurité d'un système (administrateur, compte de service
puissant). La \emph{gestion des accès à privilèges} (PAM) les place dans
un coffre, délivre leurs mots de passe à la demande, enregistre les
sessions et les renouvelle après usage.
\end{boxdef}

% ============================================================================
\section{La Revue des Droits}
\label{sec:c17-revue}
% ============================================================================

\begin{boxscenario}{Opération GAROUA. La revue des 212 comptes}
\textbf{Constatations} (rapprochement de l'annuaire et du fichier du
personnel, 15~mars)~:

\medskip
\begin{tableaucours}{L{6.6cm}c}{%
    \tableauHeader{Catégorie} &
    \tableauHeader{Comptes}
}
Comptes nominatifs de salariés présents & 148 \\
Comptes de salariés partis (orphelins) & 27 \\
Comptes partagés (agences, prestataire) & 9 \\
Comptes de service & 18 \\
Comptes sans titulaire identifiable & 10 \\
\textbf{Total} & \textbf{212} \\
\end{tableaucours}

\medskip
Deux salariés présents n'ont pas de compte (148 sur 150). Parmi les
comptes actifs, 14 ont des droits d'administrateur du domaine, dont
3 comptes orphelins et le compte partagé du prestataire. Aucun compte
distant n'a de second facteur.

\medskip
\textbf{Analyse} (confiance élevée). 46~comptes (orphelins, partagés et
sans titulaire) ne correspondent à aucune personne responsable~: ce sont
autant de portes que personne ne surveille.
\end{boxscenario}

\begin{boxobjet}{M-17 · Revue trimestrielle des comptes}
\textbf{Objectif~:} qu'aucun compte ne reste actif sans titulaire ni
besoin. \textbf{Mise en œuvre~:} extraction de l'annuaire, rapprochement
avec le fichier du personnel, validation des droits par chaque
responsable, désactivation des écarts sous huit jours.
\textbf{Vérification~:} l'auditeur refait le rapprochement et compte les
écarts restants.
\end{boxobjet}

\begin{boxmetier}{Administrateur systèmes et réseaux}
\textbf{Réduire les administrateurs.} L'administrateur ramène les comptes
d'administrateur du domaine au strict nécessaire (quelques personnes
nommées), sépare leurs comptes d'administration de leurs comptes de
travail, et active le second facteur sur tous.
\end{boxmetier}

\begin{boxmetier}{Auditeur de sécurité}
\textbf{Refaire le rapprochement.} L'auditeur ne se contente pas de la
procédure~: il extrait lui-même les comptes et les rapproche du
personnel. Le nombre d'écarts est la preuve de l'efficacité de la revue.
\end{boxmetier}

\separateur

\begin{boxresume}
\textbf{Ce qu'il faut retenir de ce chapitre~:}
\begin{itemize}
    \item \textbf{IAAA}~: identification, authentification, autorisation,
          traçabilité~; un compte partagé ruine la traçabilité.
    \item \textbf{Les mots de passe} se veulent longs, sans changement
          périodique imposé, comparés aux listes divulguées.
    \item \textbf{Le second facteur} combine des catégories différentes~;
          FIDO2 et WebAuthn résistent à l'hameçonnage.
    \item \textbf{Le SSO} réduit les mots de passe et concentre le risque
          sur le fournisseur d'identité.
    \item \textbf{Le cycle de vie} des comptes et la gestion des privilèges
          ferment les portes oubliées~; la revue des droits le vérifie.
\end{itemize}
\end{boxresume}

\begin{boxinfo}
\textbf{Points de contrôle de l'élève.} Je sais~: distinguer les quatre
étapes IAAA~; citer les trois catégories de facteurs~; énoncer la
politique de mots de passe du NIST~; expliquer pourquoi FIDO2 résiste à
l'hameçonnage~; décrire le cycle de vie d'un compte~; conduire un
rapprochement comptes-personnel.
\end{boxinfo}

\bigskip

\begin{boxexo}[title={\ding{45}\quad Exercice~17.1. IAAA \niveaubase}]
Classez~: saisir son badge~; valider un code reçu~; recevoir l'accès au
module de paie~; consulter le journal de ses connexions.
\end{boxexo}

\begin{boxexo}[title={\ding{45}\quad Exercice~17.2. Multifacteur~? \niveaubase}]
Lesquelles de ces combinaisons sont multifacteurs~: mot de passe et
question secrète~; carte et code PIN~; empreinte et visage~; mot de passe
et clé FIDO2~?
\end{boxexo}

\begin{boxexo}[title={\ding{45}\quad Exercice~17.3. Les chiffres de la revue \niveaubase}]
Vérifiez que les catégories du scénario totalisent 212 comptes, et
calculez la part des comptes sans responsable.
\end{boxexo}

\begin{boxexo}[title={\ding{45}\quad Exercice~17.4. La politique de mots de passe \niveaumoyen}]
La politique actuelle de la BANQUE SANAGA impose 8~caractères, quatre
types de caractères et un changement tous les 30~jours. Réécrivez-la selon
le NIST SP~800-63B et justifiez chaque changement.
\end{boxexo}

\begin{boxexo}[title={\ding{45}\quad Exercice~17.5. Le relais \niveaumoyen}]
Décrivez comment une fausse page peut relayer en temps réel un code TOTP,
et pourquoi une clé FIDO2 l'empêche.
\end{boxexo}

\begin{boxexo}[title={\ding{45}\quad Exercice~17.6. Le compte du prestataire \niveaumoyen}]
Proposez l'organisation des accès de SAHEL INFOGÉRANCE~: comptes,
facteurs, droits, horaires, journalisation.
\end{boxexo}

\begin{boxexo}[title={\ding{45}\quad Exercice~17.7. Comptes de service \niveaumoyen}]
Les 18~comptes de service de LiZbiZ ont des mots de passe qui ne changent
jamais et des droits d'administrateur. Quels risques~? Quelles mesures~?
\end{boxexo}

\begin{boxexo}[title={\ding{45}\quad Exercice~17.8. SSO \niveaumoyen}]
Quels avantages et quels risques l'authentification unique apporte-t-elle
à LiZbiZ~? Quelles protections exige le fournisseur d'identité~?
\end{boxexo}

\begin{boxexo}[title={\ding{45}\quad Exercice~17.9. Le plan de remédiation \niveauexpert}]
À partir de la revue des 212~comptes, rédigez un plan de remédiation en
trente jours, par priorité, avec les responsables.
\end{boxexo}

\begin{boxexo}[title={\ding{45}\quad Exercice~17.10. Un script de rapprochement \niveauexpert}]
Écrivez, en Python, un programme qui lit un export de l'annuaire et le
fichier du personnel (deux fichiers CSV) et produit la liste des comptes
orphelins, partagés et sans titulaire.
\end{boxexo}

\bigskip

\begin{boxinfo}
\textbf{Pour aller plus loin.}
\begin{itemize}
    \item NIST SP~800-63B~\cite{nist80063b}.
    \item W3C, \emph{Web Authentication}~\cite{w3c2019webauthn}~; SAML~2.0
          et OpenID Connect~\cite{oasis2005saml,oidc2014}.
\end{itemize}
\end{boxinfo}

\begin{boxinfo}
\textbf{Transition.} Les comptes donnent accès à des systèmes~; encore
faut-il que ces systèmes soient solides. Le Chapitre~\ref{chap:systemes}
traite de leur durcissement et de leur maintien en condition de sécurité.
\end{boxinfo}
